% Appendix B: Methodology for in-the-loop FEM
% Include this in your main .tex: % Appendix B: Methodology for in-the-loop FEM
% Include this in your main .tex: % Appendix B: Methodology for in-the-loop FEM
% Include this in your main .tex: % Appendix B: Methodology for in-the-loop FEM
% Include this in your main .tex: \input{docs/appendix_B_fem} or \include{...}

\section{Methodology for in-the-loop FEM}\label{sec:appendix_B}

FEM calculations must be optimized for speed in order to be computed in the optimization loop. We detail below a number of choices that were made to implement, make efficient, and make robust the in-the-loop FEM calculations presented in this work.

\paragraph{Domain and Geometry.}
The coil winding pack is represented as a solid with rectangular cross-section swept along the coil centerline. The mesh is a tetrahedralization of this volume, generated either by Gmsh or a pre-built mesh from ParaStell or Gmsh.

\paragraph{Mesh Deformation Strategy.}
The mesh topology is generated once and cached. On subsequent evaluations, node positions are deformed to follow the updated coil geometry:
\begin{enumerate}
  \item For each mesh node $\mathbf{x}$, find the nearest coil centerline point $\boldsymbol{\gamma}_{i}$ and the associated coil frame $\{\mathbf{p}_i, \mathbf{q}_i, \mathbf{t}_i\}$.
  \item Compute the local offset in the coil frame:
    \begin{equation}
      \delta_p = (\mathbf{x} - \boldsymbol{\gamma}_i) \cdot \mathbf{p}_i,\quad
      \delta_q = (\mathbf{x} - \boldsymbol{\gamma}_i) \cdot \mathbf{q}_i,\quad
      \delta_t = (\mathbf{x} - \boldsymbol{\gamma}_i) \cdot \mathbf{t}_i.
    \end{equation}
  \item When the coil deforms, reconstruct the new node position:
    \begin{equation}
      \mathbf{x}_{\text{new}} = \boldsymbol{\gamma}_i + \delta_p \, \mathbf{p}_i + \delta_q \, \mathbf{q}_i + \delta_t \, \mathbf{t}_i.
    \end{equation}
\end{enumerate}

The coil frame uses the centroid frame: $\mathbf{C} = \operatorname{mean}(\boldsymbol{\gamma})$, $\mathbf{w} = \boldsymbol{\gamma} - \mathbf{C}$, $\mathbf{p} = (\mathbf{w} - (\mathbf{w} \cdot \mathbf{t})\mathbf{t}) / |\cdots|$, $\mathbf{q} = \mathbf{t} \times \mathbf{p}$.

\paragraph{Current Density Model.}
The body force density is $\bm f = \bm J \times \bm B$, which requires a model for the current in the coil and the magnetic field and must be evaluated on the appropriate quadrature points of the weak formulation described below.
The current density is assumed uniform across the winding-pack cross-section:
\begin{equation}
  \mathbf{J} = \frac{I}{A} \, \mathbf{t},
\end{equation}
where $I$ is the coil current, $A$ is the cross-section area, and $\mathbf{t}$ is the unit tangent at the nearest centerline point. Each mesh node (or quadrature point) is assigned to the nearest coil centerline point via a KD-tree lookup.

\paragraph{Magnetic Field model.}
The Biot--Savart law diverges near the coil centerline, so a filamentary model is inadequate for points inside the conductor. StellCoilBench uses the Landreman et al.\ (2025) regularized internal field model for rectangular cross-sections. The total field at a point inside a coil is:
\begin{equation}
  \mathbf{B} = \mathbf{B}_{\text{reg}} + \mathbf{B}_0 + \mathbf{B}_{\kappa} + \mathbf{B}_b.
\end{equation}
This is a fast approximation for the self-field that significantly increases the capacity to compute the Lorentz body force in-the-loop.
The mutual field is computed via the Biot-Savart law assuming filamentary models for the other coils, since this is adequate outside the coil where the calculations are taking place (assuming other coils do not come too close).

\paragraph{Linear Elasticity.}
The linear structural problem to solve is: find the displacement $\mathbf{u} : \Omega \to \mathbb{R}^3$ such that
\begin{equation}
  -\nabla \cdot \boldsymbol{\sigma} = \mathbf{f}
  \quad \text{in } \Omega,
  \qquad
  \mathbf{u} = \mathbf{0}
  \quad \text{on } \Gamma_\Omega,
\end{equation}
where $\boldsymbol{\sigma}$ is the Cauchy stress tensor for a linear elastic, isotropic material:
\begin{equation}
  \boldsymbol{\sigma} = \lambda \, (\nabla \cdot \mathbf{u}) \, \mathbf{I} + 2\mu \, \boldsymbol{\varepsilon},
  \qquad
  \boldsymbol{\varepsilon} = \frac{1}{2}\bigl(\nabla \mathbf{u} + (\nabla \mathbf{u})^T\bigr).
\end{equation}
The Lam\'e Parameters are:
\begin{equation}
  \lambda = \frac{E \nu}{(1+\nu)(1-2\nu)},
  \qquad
  \mu = \frac{E}{2(1+\nu)},
\end{equation}
where $E$ is Young's modulus [Pa] and $\nu$ is Poisson ratio. The defaults are $E = 100\,\text{GPa}$, $\nu = 0.3$, which are typical for REBCO and similar materials~\cite{srivastava2024modelling}.
Simple ``fixed support'' boundary conditions are used; nodes whose $z$-coordinate lies in the bottom fraction $\zeta = 0.15$ of the global mesh $z$-range are pinned ($\mathbf{u} = \mathbf{0}$). This models coil support stands and removes rigid-body modes:
\begin{equation}
  \Gamma_\Omega = \bigl\{ \mathbf{x} \in \partial\Omega : z(\mathbf{x}) \leq z_{\min} + 0.15(z_{\max} - z_{\min}) \bigr\}.
\end{equation}
More complex boundary conditions consistent with realistic support structures can easily be implemented in the code in future work.

\paragraph{Weak Form and Discretization.}
Using linear tetrahedral elements and a standard Galerkin discretization:
\begin{equation}
  \int_\Omega \bigl( \lambda \, (\nabla \cdot \mathbf{u}) \, (\nabla \cdot \mathbf{v}) + 2\mu \, \boldsymbol{\varepsilon}(\mathbf{u}) : \boldsymbol{\varepsilon}(\mathbf{v}) \bigr) \, dV
  = \int_\Omega \mathbf{f} \cdot \mathbf{v} \, dV
  \quad \forall \mathbf{v} \in V_h.
\end{equation}
The quadrature degree for the load integral is configurable (see Reduced Quadrature below).
The implementation uses scikit-fem or DOLFINx with bilinear and linear forms assembled and the linear system solved with Dirichlet DOFs eliminated via condense.

\paragraph{Von Mises Stress.}
For each tetrahedral element, the displacement gradient $\nabla \mathbf{u}$ is computed from nodal values and the element Jacobian. The strain tensor is
\begin{equation}
  \boldsymbol{\varepsilon} = \frac{1}{2}\bigl(\nabla \mathbf{u} + (\nabla \mathbf{u})^T\bigr).
\end{equation}
The Cauchy stress is
\begin{equation}
  \boldsymbol{\sigma} = \lambda \, \operatorname{tr}(\boldsymbol{\varepsilon}) \, \mathbf{I} + 2\mu \, \boldsymbol{\varepsilon}.
\end{equation}
The deviatoric stress is $\mathbf{s} = \boldsymbol{\sigma} - \frac{1}{3}\operatorname{tr}(\boldsymbol{\sigma}) \, \mathbf{I}$. The Von Mises stress is
\begin{equation}
  \sigma_{\text{vm}}
  = \sqrt{\frac{3}{2} \sum_{i,j} s_{ij}^2}.
\end{equation}
This is computed per element (constant within each linear tetrahedron).
The FEM yields $\sigma_{\text{vm}}^{(e)}$ per element $e$, so that the objective is chosen to be:
\begin{equation}
  \max(0, \max_e \sigma_{\text{vm}}^{(e)} - \lambda_{\sigma_{vm}})^2.
\end{equation}
The objective returns values in GPa to improve optimizer numerics.

\paragraph{Optimization tricks.}
We use a number of optimization tricks to prevent undesirable behavior.
The first is a guard against computing FEM when coil-coil distances become very small, since the mutual field body force calculations will become nearly singular. When $d_{cc} < \alpha_\text{safety}\lambda_{d_{cc}}$ it skips the FEM and returns a finite penalty (e.g., 10\,GPa) with zero gradient.
The second trick is a short-circuit; when $\max_e \sigma_{\text{vm}}^{(e)} \leq \lambda_{\sigma_{vm}}$, it skips the expensive gradient calculation and returns $\nabla J = \mathbf{0}$.
Gradients are computed by forward finite differences.
Third, the stiffness matrix $\mathbf{K}$ and fixed-DOF list can be cached at the baseline; $\mathbf{K}$ changes by $O(\varepsilon)$ per perturbation, which is negligible for FD accuracy. When \texttt{structural\_use\_cached\_K} is true, $\mathbf{K}$ is reused for all FD perturbations in $\nabla J$, yielding up to $\sim$40\,\% speedup on $\nabla J$ for adaptive meshes at the cost of gradient accuracy; the default is false.
Fourth, only coil frames and $\mathbf{B}_{\text{reg}}$ for coils whose DOFs are perturbed are recomputed, reducing cost.
Fifth, the objective can be evaluated every $N$ iterations and cached in between to reduce cost.
Both DOLFINx and scikit-fem evaluate J$\times$B at quadrature points; the backend choice affects per-solve speed (see Dual Backends below).

\paragraph{MPI Parallelization.}
When running with MPI ($P > 1$ processes), the structural gradient $\nabla J$ is computed in parallel: rank~0 broadcasts the baseline state; each rank computes a subset of the $n$ finite-difference perturbations (indices $r, r+P, r+2P, \ldots$); partial gradients are combined via \texttt{Allreduce}. Workers (ranks $1 \ldots P{-}1$) block in a worker loop until rank~0 signals a gradient computation. Expected speedup is up to $\sim P\times$ on the FD loop; overhead from broadcast and Allreduce is small for typical $n \sim 100$--200. Example: \texttt{mpirun -n 4 stellcoilbench submit-case cases/basic\_LandremanPaulQA\_structural.yaml} uses 4 processes for parallel $\nabla J$.

\paragraph{Reduced Quadrature.}
The weak-form RHS uses a quadrature element with configurable degree $q \in \{1,2\}$. For $q=1$ (centroid-only), J$\times$B evaluations are reduced, giving $\sim$30--50\,\% speedup in body-force assembly with $\sim$2\,\% relative error in $J$ for typical coils. Configurable via \texttt{structural\_quadrature\_degree}.

\paragraph{Adaptive Mesh.}
When \texttt{structural\_mesh\_resolution\_coarse} and \texttt{structural\_mesh\_resolution\_fine} are set, optimization starts with the coarse mesh and refines to the fine mesh when $\sigma_{\text{vm}} \geq \eta \times \lambda_{\sigma_{\text{vm}}}$ ($\eta = 0.5$ by default). Refinement is one-way (coarse $\to$ fine) and triggered from the optimizer callback; subsequent iterations use the finer mesh. Early iterations are 2--6$\times$ faster.

\paragraph{Dual Backends.}
DOLFINx (default) and scikit-fem are supported. DOLFINx is typically 2.5--3$\times$ faster per solve; scikit-fem is used when DOLFINx is unavailable or when \texttt{structural\_backend: skfem} is set. Both evaluate J$\times$B at quadrature points and use P1 tetrahedra.

\begin{table*}[t]
  \centering
  \caption{Key implementation parameters.}
  \begin{tabular}{lll}
    \toprule
    Parameter & Default & Description \\
    \midrule
    \texttt{width}, \texttt{height} & 0.05\,m & Cross-section dimensions \\
    \texttt{E} & 100\,GPa & Young's modulus \\
    \texttt{nu} & 0.3 & Poisson ratio \\
    \texttt{mesh\_resolution} & 0.16\,m & Gmsh element size \\
    \texttt{stress\_metric} & $\max_e \sigma_{\text{vm}}^{(e)}$ & Scalar metric \\
    \texttt{fd\_step} & $10^{-5}$ & Finite-difference step \\
    \texttt{eval\_interval} & 1 & Evaluate every $N$ iterations \\
    \texttt{BC\_Z\_FRACTION} & 0.15 & Fraction of $z$-range pinned \\
    \texttt{structural\_quadrature\_degree} & 2 & Quadrature degree for J$\times$B (1=centroid, 2=full) \\
    \texttt{structural\_mesh\_resolution\_coarse} & 0.16\,m & Coarse mesh for adaptive strategy \\
    \texttt{structural\_mesh\_resolution\_fine} & 0.08\,m & Fine mesh after refinement \\
    \texttt{structural\_refine\_stress\_ratio} & 0.5 & Refine when $\sigma_{\text{vm}} \geq \eta \times$ threshold \\
    \texttt{structural\_use\_cached\_K} & false & Cache stiffness $\mathbf{K}$ for $\nabla J$ speedup \\
    \texttt{structural\_backend} & dolfinx & FEM backend: dolfinx or skfem \\
    \bottomrule
  \end{tabular}
  \label{tab:fem_summary}
\end{table*}
 or \include{...}

\section{Methodology for in-the-loop FEM}\label{sec:appendix_B}

FEM calculations must be optimized for speed in order to be computed in the optimization loop. We detail below a number of choices that were made to implement, make efficient, and make robust the in-the-loop FEM calculations presented in this work.

\paragraph{Domain and Geometry.}
The coil winding pack is represented as a solid with rectangular cross-section swept along the coil centerline. The mesh is a tetrahedralization of this volume, generated either by Gmsh or a pre-built mesh from ParaStell or Gmsh.

\paragraph{Mesh Deformation Strategy.}
The mesh topology is generated once and cached. On subsequent evaluations, node positions are deformed to follow the updated coil geometry:
\begin{enumerate}
  \item For each mesh node $\mathbf{x}$, find the nearest coil centerline point $\boldsymbol{\gamma}_{i}$ and the associated coil frame $\{\mathbf{p}_i, \mathbf{q}_i, \mathbf{t}_i\}$.
  \item Compute the local offset in the coil frame:
    \begin{equation}
      \delta_p = (\mathbf{x} - \boldsymbol{\gamma}_i) \cdot \mathbf{p}_i,\quad
      \delta_q = (\mathbf{x} - \boldsymbol{\gamma}_i) \cdot \mathbf{q}_i,\quad
      \delta_t = (\mathbf{x} - \boldsymbol{\gamma}_i) \cdot \mathbf{t}_i.
    \end{equation}
  \item When the coil deforms, reconstruct the new node position:
    \begin{equation}
      \mathbf{x}_{\text{new}} = \boldsymbol{\gamma}_i + \delta_p \, \mathbf{p}_i + \delta_q \, \mathbf{q}_i + \delta_t \, \mathbf{t}_i.
    \end{equation}
\end{enumerate}

The coil frame uses the centroid frame: $\mathbf{C} = \operatorname{mean}(\boldsymbol{\gamma})$, $\mathbf{w} = \boldsymbol{\gamma} - \mathbf{C}$, $\mathbf{p} = (\mathbf{w} - (\mathbf{w} \cdot \mathbf{t})\mathbf{t}) / |\cdots|$, $\mathbf{q} = \mathbf{t} \times \mathbf{p}$.

\paragraph{Current Density Model.}
The body force density is $\bm f = \bm J \times \bm B$, which requires a model for the current in the coil and the magnetic field and must be evaluated on the appropriate quadrature points of the weak formulation described below.
The current density is assumed uniform across the winding-pack cross-section:
\begin{equation}
  \mathbf{J} = \frac{I}{A} \, \mathbf{t},
\end{equation}
where $I$ is the coil current, $A$ is the cross-section area, and $\mathbf{t}$ is the unit tangent at the nearest centerline point. Each mesh node (or quadrature point) is assigned to the nearest coil centerline point via a KD-tree lookup.

\paragraph{Magnetic Field model.}
The Biot--Savart law diverges near the coil centerline, so a filamentary model is inadequate for points inside the conductor. StellCoilBench uses the Landreman et al.\ (2025) regularized internal field model for rectangular cross-sections. The total field at a point inside a coil is:
\begin{equation}
  \mathbf{B} = \mathbf{B}_{\text{reg}} + \mathbf{B}_0 + \mathbf{B}_{\kappa} + \mathbf{B}_b.
\end{equation}
This is a fast approximation for the self-field that significantly increases the capacity to compute the Lorentz body force in-the-loop.
The mutual field is computed via the Biot-Savart law assuming filamentary models for the other coils, since this is adequate outside the coil where the calculations are taking place (assuming other coils do not come too close).

\paragraph{Linear Elasticity.}
The linear structural problem to solve is: find the displacement $\mathbf{u} : \Omega \to \mathbb{R}^3$ such that
\begin{equation}
  -\nabla \cdot \boldsymbol{\sigma} = \mathbf{f}
  \quad \text{in } \Omega,
  \qquad
  \mathbf{u} = \mathbf{0}
  \quad \text{on } \Gamma_\Omega,
\end{equation}
where $\boldsymbol{\sigma}$ is the Cauchy stress tensor for a linear elastic, isotropic material:
\begin{equation}
  \boldsymbol{\sigma} = \lambda \, (\nabla \cdot \mathbf{u}) \, \mathbf{I} + 2\mu \, \boldsymbol{\varepsilon},
  \qquad
  \boldsymbol{\varepsilon} = \frac{1}{2}\bigl(\nabla \mathbf{u} + (\nabla \mathbf{u})^T\bigr).
\end{equation}
The Lam\'e Parameters are:
\begin{equation}
  \lambda = \frac{E \nu}{(1+\nu)(1-2\nu)},
  \qquad
  \mu = \frac{E}{2(1+\nu)},
\end{equation}
where $E$ is Young's modulus [Pa] and $\nu$ is Poisson ratio. The defaults are $E = 100\,\text{GPa}$, $\nu = 0.3$, which are typical for REBCO and similar materials~\cite{srivastava2024modelling}.
Simple ``fixed support'' boundary conditions are used; nodes whose $z$-coordinate lies in the bottom fraction $\zeta = 0.15$ of the global mesh $z$-range are pinned ($\mathbf{u} = \mathbf{0}$). This models coil support stands and removes rigid-body modes:
\begin{equation}
  \Gamma_\Omega = \bigl\{ \mathbf{x} \in \partial\Omega : z(\mathbf{x}) \leq z_{\min} + 0.15(z_{\max} - z_{\min}) \bigr\}.
\end{equation}
More complex boundary conditions consistent with realistic support structures can easily be implemented in the code in future work.

\paragraph{Weak Form and Discretization.}
Using linear tetrahedral elements and a standard Galerkin discretization:
\begin{equation}
  \int_\Omega \bigl( \lambda \, (\nabla \cdot \mathbf{u}) \, (\nabla \cdot \mathbf{v}) + 2\mu \, \boldsymbol{\varepsilon}(\mathbf{u}) : \boldsymbol{\varepsilon}(\mathbf{v}) \bigr) \, dV
  = \int_\Omega \mathbf{f} \cdot \mathbf{v} \, dV
  \quad \forall \mathbf{v} \in V_h.
\end{equation}
The quadrature degree for the load integral is configurable (see Reduced Quadrature below).
The implementation uses scikit-fem or DOLFINx with bilinear and linear forms assembled and the linear system solved with Dirichlet DOFs eliminated via condense.

\paragraph{Von Mises Stress.}
For each tetrahedral element, the displacement gradient $\nabla \mathbf{u}$ is computed from nodal values and the element Jacobian. The strain tensor is
\begin{equation}
  \boldsymbol{\varepsilon} = \frac{1}{2}\bigl(\nabla \mathbf{u} + (\nabla \mathbf{u})^T\bigr).
\end{equation}
The Cauchy stress is
\begin{equation}
  \boldsymbol{\sigma} = \lambda \, \operatorname{tr}(\boldsymbol{\varepsilon}) \, \mathbf{I} + 2\mu \, \boldsymbol{\varepsilon}.
\end{equation}
The deviatoric stress is $\mathbf{s} = \boldsymbol{\sigma} - \frac{1}{3}\operatorname{tr}(\boldsymbol{\sigma}) \, \mathbf{I}$. The Von Mises stress is
\begin{equation}
  \sigma_{\text{vm}}
  = \sqrt{\frac{3}{2} \sum_{i,j} s_{ij}^2}.
\end{equation}
This is computed per element (constant within each linear tetrahedron).
The FEM yields $\sigma_{\text{vm}}^{(e)}$ per element $e$, so that the objective is chosen to be:
\begin{equation}
  \max(0, \max_e \sigma_{\text{vm}}^{(e)} - \lambda_{\sigma_{vm}})^2.
\end{equation}
The objective returns values in GPa to improve optimizer numerics.

\paragraph{Optimization tricks.}
We use a number of optimization tricks to prevent undesirable behavior.
The first is a guard against computing FEM when coil-coil distances become very small, since the mutual field body force calculations will become nearly singular. When $d_{cc} < \alpha_\text{safety}\lambda_{d_{cc}}$ it skips the FEM and returns a finite penalty (e.g., 10\,GPa) with zero gradient.
The second trick is a short-circuit; when $\max_e \sigma_{\text{vm}}^{(e)} \leq \lambda_{\sigma_{vm}}$, it skips the expensive gradient calculation and returns $\nabla J = \mathbf{0}$.
Gradients are computed by forward finite differences.
Third, the stiffness matrix $\mathbf{K}$ and fixed-DOF list can be cached at the baseline; $\mathbf{K}$ changes by $O(\varepsilon)$ per perturbation, which is negligible for FD accuracy. When \texttt{structural\_use\_cached\_K} is true, $\mathbf{K}$ is reused for all FD perturbations in $\nabla J$, yielding up to $\sim$40\,\% speedup on $\nabla J$ for adaptive meshes at the cost of gradient accuracy; the default is false.
Fourth, only coil frames and $\mathbf{B}_{\text{reg}}$ for coils whose DOFs are perturbed are recomputed, reducing cost.
Fifth, the objective can be evaluated every $N$ iterations and cached in between to reduce cost.
Both DOLFINx and scikit-fem evaluate J$\times$B at quadrature points; the backend choice affects per-solve speed (see Dual Backends below).

\paragraph{MPI Parallelization.}
When running with MPI ($P > 1$ processes), the structural gradient $\nabla J$ is computed in parallel: rank~0 broadcasts the baseline state; each rank computes a subset of the $n$ finite-difference perturbations (indices $r, r+P, r+2P, \ldots$); partial gradients are combined via \texttt{Allreduce}. Workers (ranks $1 \ldots P{-}1$) block in a worker loop until rank~0 signals a gradient computation. Expected speedup is up to $\sim P\times$ on the FD loop; overhead from broadcast and Allreduce is small for typical $n \sim 100$--200. Example: \texttt{mpirun -n 4 stellcoilbench submit-case cases/basic\_LandremanPaulQA\_structural.yaml} uses 4 processes for parallel $\nabla J$.

\paragraph{Reduced Quadrature.}
The weak-form RHS uses a quadrature element with configurable degree $q \in \{1,2\}$. For $q=1$ (centroid-only), J$\times$B evaluations are reduced, giving $\sim$30--50\,\% speedup in body-force assembly with $\sim$2\,\% relative error in $J$ for typical coils. Configurable via \texttt{structural\_quadrature\_degree}.

\paragraph{Adaptive Mesh.}
When \texttt{structural\_mesh\_resolution\_coarse} and \texttt{structural\_mesh\_resolution\_fine} are set, optimization starts with the coarse mesh and refines to the fine mesh when $\sigma_{\text{vm}} \geq \eta \times \lambda_{\sigma_{\text{vm}}}$ ($\eta = 0.5$ by default). Refinement is one-way (coarse $\to$ fine) and triggered from the optimizer callback; subsequent iterations use the finer mesh. Early iterations are 2--6$\times$ faster.

\paragraph{Dual Backends.}
DOLFINx (default) and scikit-fem are supported. DOLFINx is typically 2.5--3$\times$ faster per solve; scikit-fem is used when DOLFINx is unavailable or when \texttt{structural\_backend: skfem} is set. Both evaluate J$\times$B at quadrature points and use P1 tetrahedra.

\begin{table*}[t]
  \centering
  \caption{Key implementation parameters.}
  \begin{tabular}{lll}
    \toprule
    Parameter & Default & Description \\
    \midrule
    \texttt{width}, \texttt{height} & 0.05\,m & Cross-section dimensions \\
    \texttt{E} & 100\,GPa & Young's modulus \\
    \texttt{nu} & 0.3 & Poisson ratio \\
    \texttt{mesh\_resolution} & 0.16\,m & Gmsh element size \\
    \texttt{stress\_metric} & $\max_e \sigma_{\text{vm}}^{(e)}$ & Scalar metric \\
    \texttt{fd\_step} & $10^{-5}$ & Finite-difference step \\
    \texttt{eval\_interval} & 1 & Evaluate every $N$ iterations \\
    \texttt{BC\_Z\_FRACTION} & 0.15 & Fraction of $z$-range pinned \\
    \texttt{structural\_quadrature\_degree} & 2 & Quadrature degree for J$\times$B (1=centroid, 2=full) \\
    \texttt{structural\_mesh\_resolution\_coarse} & 0.16\,m & Coarse mesh for adaptive strategy \\
    \texttt{structural\_mesh\_resolution\_fine} & 0.08\,m & Fine mesh after refinement \\
    \texttt{structural\_refine\_stress\_ratio} & 0.5 & Refine when $\sigma_{\text{vm}} \geq \eta \times$ threshold \\
    \texttt{structural\_use\_cached\_K} & false & Cache stiffness $\mathbf{K}$ for $\nabla J$ speedup \\
    \texttt{structural\_backend} & dolfinx & FEM backend: dolfinx or skfem \\
    \bottomrule
  \end{tabular}
  \label{tab:fem_summary}
\end{table*}
 or \include{...}

\section{Methodology for in-the-loop FEM}\label{sec:appendix_B}

FEM calculations must be optimized for speed in order to be computed in the optimization loop. We detail below a number of choices that were made to implement, make efficient, and make robust the in-the-loop FEM calculations presented in this work.

\paragraph{Domain and Geometry.}
The coil winding pack is represented as a solid with rectangular cross-section swept along the coil centerline. The mesh is a tetrahedralization of this volume, generated either by Gmsh or a pre-built mesh from ParaStell or Gmsh.

\paragraph{Mesh Deformation Strategy.}
The mesh topology is generated once and cached. On subsequent evaluations, node positions are deformed to follow the updated coil geometry:
\begin{enumerate}
  \item For each mesh node $\mathbf{x}$, find the nearest coil centerline point $\boldsymbol{\gamma}_{i}$ and the associated coil frame $\{\mathbf{p}_i, \mathbf{q}_i, \mathbf{t}_i\}$.
  \item Compute the local offset in the coil frame:
    \begin{equation}
      \delta_p = (\mathbf{x} - \boldsymbol{\gamma}_i) \cdot \mathbf{p}_i,\quad
      \delta_q = (\mathbf{x} - \boldsymbol{\gamma}_i) \cdot \mathbf{q}_i,\quad
      \delta_t = (\mathbf{x} - \boldsymbol{\gamma}_i) \cdot \mathbf{t}_i.
    \end{equation}
  \item When the coil deforms, reconstruct the new node position:
    \begin{equation}
      \mathbf{x}_{\text{new}} = \boldsymbol{\gamma}_i + \delta_p \, \mathbf{p}_i + \delta_q \, \mathbf{q}_i + \delta_t \, \mathbf{t}_i.
    \end{equation}
\end{enumerate}

The coil frame uses the centroid frame: $\mathbf{C} = \operatorname{mean}(\boldsymbol{\gamma})$, $\mathbf{w} = \boldsymbol{\gamma} - \mathbf{C}$, $\mathbf{p} = (\mathbf{w} - (\mathbf{w} \cdot \mathbf{t})\mathbf{t}) / |\cdots|$, $\mathbf{q} = \mathbf{t} \times \mathbf{p}$.

\paragraph{Current Density Model.}
The body force density is $\bm f = \bm J \times \bm B$, which requires a model for the current in the coil and the magnetic field and must be evaluated on the appropriate quadrature points of the weak formulation described below.
The current density is assumed uniform across the winding-pack cross-section:
\begin{equation}
  \mathbf{J} = \frac{I}{A} \, \mathbf{t},
\end{equation}
where $I$ is the coil current, $A$ is the cross-section area, and $\mathbf{t}$ is the unit tangent at the nearest centerline point. Each mesh node (or quadrature point) is assigned to the nearest coil centerline point via a KD-tree lookup.

\paragraph{Magnetic Field model.}
The Biot--Savart law diverges near the coil centerline, so a filamentary model is inadequate for points inside the conductor. StellCoilBench uses the Landreman et al.\ (2025) regularized internal field model for rectangular cross-sections. The total field at a point inside a coil is:
\begin{equation}
  \mathbf{B} = \mathbf{B}_{\text{reg}} + \mathbf{B}_0 + \mathbf{B}_{\kappa} + \mathbf{B}_b.
\end{equation}
This is a fast approximation for the self-field that significantly increases the capacity to compute the Lorentz body force in-the-loop.
The mutual field is computed via the Biot-Savart law assuming filamentary models for the other coils, since this is adequate outside the coil where the calculations are taking place (assuming other coils do not come too close).

\paragraph{Linear Elasticity.}
The linear structural problem to solve is: find the displacement $\mathbf{u} : \Omega \to \mathbb{R}^3$ such that
\begin{equation}
  -\nabla \cdot \boldsymbol{\sigma} = \mathbf{f}
  \quad \text{in } \Omega,
  \qquad
  \mathbf{u} = \mathbf{0}
  \quad \text{on } \Gamma_\Omega,
\end{equation}
where $\boldsymbol{\sigma}$ is the Cauchy stress tensor for a linear elastic, isotropic material:
\begin{equation}
  \boldsymbol{\sigma} = \lambda \, (\nabla \cdot \mathbf{u}) \, \mathbf{I} + 2\mu \, \boldsymbol{\varepsilon},
  \qquad
  \boldsymbol{\varepsilon} = \frac{1}{2}\bigl(\nabla \mathbf{u} + (\nabla \mathbf{u})^T\bigr).
\end{equation}
The Lam\'e Parameters are:
\begin{equation}
  \lambda = \frac{E \nu}{(1+\nu)(1-2\nu)},
  \qquad
  \mu = \frac{E}{2(1+\nu)},
\end{equation}
where $E$ is Young's modulus [Pa] and $\nu$ is Poisson ratio. The defaults are $E = 100\,\text{GPa}$, $\nu = 0.3$, which are typical for REBCO and similar materials~\cite{srivastava2024modelling}.
Simple ``fixed support'' boundary conditions are used; nodes whose $z$-coordinate lies in the bottom fraction $\zeta = 0.15$ of the global mesh $z$-range are pinned ($\mathbf{u} = \mathbf{0}$). This models coil support stands and removes rigid-body modes:
\begin{equation}
  \Gamma_\Omega = \bigl\{ \mathbf{x} \in \partial\Omega : z(\mathbf{x}) \leq z_{\min} + 0.15(z_{\max} - z_{\min}) \bigr\}.
\end{equation}
More complex boundary conditions consistent with realistic support structures can easily be implemented in the code in future work.

\paragraph{Weak Form and Discretization.}
Using linear tetrahedral elements and a standard Galerkin discretization:
\begin{equation}
  \int_\Omega \bigl( \lambda \, (\nabla \cdot \mathbf{u}) \, (\nabla \cdot \mathbf{v}) + 2\mu \, \boldsymbol{\varepsilon}(\mathbf{u}) : \boldsymbol{\varepsilon}(\mathbf{v}) \bigr) \, dV
  = \int_\Omega \mathbf{f} \cdot \mathbf{v} \, dV
  \quad \forall \mathbf{v} \in V_h.
\end{equation}
The quadrature degree for the load integral is configurable (see Reduced Quadrature below).
The implementation uses scikit-fem or DOLFINx with bilinear and linear forms assembled and the linear system solved with Dirichlet DOFs eliminated via condense.

\paragraph{Von Mises Stress.}
For each tetrahedral element, the displacement gradient $\nabla \mathbf{u}$ is computed from nodal values and the element Jacobian. The strain tensor is
\begin{equation}
  \boldsymbol{\varepsilon} = \frac{1}{2}\bigl(\nabla \mathbf{u} + (\nabla \mathbf{u})^T\bigr).
\end{equation}
The Cauchy stress is
\begin{equation}
  \boldsymbol{\sigma} = \lambda \, \operatorname{tr}(\boldsymbol{\varepsilon}) \, \mathbf{I} + 2\mu \, \boldsymbol{\varepsilon}.
\end{equation}
The deviatoric stress is $\mathbf{s} = \boldsymbol{\sigma} - \frac{1}{3}\operatorname{tr}(\boldsymbol{\sigma}) \, \mathbf{I}$. The Von Mises stress is
\begin{equation}
  \sigma_{\text{vm}}
  = \sqrt{\frac{3}{2} \sum_{i,j} s_{ij}^2}.
\end{equation}
This is computed per element (constant within each linear tetrahedron).
The FEM yields $\sigma_{\text{vm}}^{(e)}$ per element $e$, so that the objective is chosen to be:
\begin{equation}
  \max(0, \max_e \sigma_{\text{vm}}^{(e)} - \lambda_{\sigma_{vm}})^2.
\end{equation}
The objective returns values in GPa to improve optimizer numerics.

\paragraph{Optimization tricks.}
We use a number of optimization tricks to prevent undesirable behavior.
The first is a guard against computing FEM when coil-coil distances become very small, since the mutual field body force calculations will become nearly singular. When $d_{cc} < \alpha_\text{safety}\lambda_{d_{cc}}$ it skips the FEM and returns a finite penalty (e.g., 10\,GPa) with zero gradient.
The second trick is a short-circuit; when $\max_e \sigma_{\text{vm}}^{(e)} \leq \lambda_{\sigma_{vm}}$, it skips the expensive gradient calculation and returns $\nabla J = \mathbf{0}$.
Gradients are computed by forward finite differences.
Third, the stiffness matrix $\mathbf{K}$ and fixed-DOF list can be cached at the baseline; $\mathbf{K}$ changes by $O(\varepsilon)$ per perturbation, which is negligible for FD accuracy. When \texttt{structural\_use\_cached\_K} is true, $\mathbf{K}$ is reused for all FD perturbations in $\nabla J$, yielding up to $\sim$40\,\% speedup on $\nabla J$ for adaptive meshes at the cost of gradient accuracy; the default is false.
Fourth, only coil frames and $\mathbf{B}_{\text{reg}}$ for coils whose DOFs are perturbed are recomputed, reducing cost.
Fifth, the objective can be evaluated every $N$ iterations and cached in between to reduce cost.
Both DOLFINx and scikit-fem evaluate J$\times$B at quadrature points; the backend choice affects per-solve speed (see Dual Backends below).

\paragraph{MPI Parallelization.}
When running with MPI ($P > 1$ processes), the structural gradient $\nabla J$ is computed in parallel: rank~0 broadcasts the baseline state; each rank computes a subset of the $n$ finite-difference perturbations (indices $r, r+P, r+2P, \ldots$); partial gradients are combined via \texttt{Allreduce}. Workers (ranks $1 \ldots P{-}1$) block in a worker loop until rank~0 signals a gradient computation. Expected speedup is up to $\sim P\times$ on the FD loop; overhead from broadcast and Allreduce is small for typical $n \sim 100$--200. Example: \texttt{mpirun -n 4 stellcoilbench submit-case cases/basic\_LandremanPaulQA\_structural.yaml} uses 4 processes for parallel $\nabla J$.

\paragraph{Reduced Quadrature.}
The weak-form RHS uses a quadrature element with configurable degree $q \in \{1,2\}$. For $q=1$ (centroid-only), J$\times$B evaluations are reduced, giving $\sim$30--50\,\% speedup in body-force assembly with $\sim$2\,\% relative error in $J$ for typical coils. Configurable via \texttt{structural\_quadrature\_degree}.

\paragraph{Adaptive Mesh.}
When \texttt{structural\_mesh\_resolution\_coarse} and \texttt{structural\_mesh\_resolution\_fine} are set, optimization starts with the coarse mesh and refines to the fine mesh when $\sigma_{\text{vm}} \geq \eta \times \lambda_{\sigma_{\text{vm}}}$ ($\eta = 0.5$ by default). Refinement is one-way (coarse $\to$ fine) and triggered from the optimizer callback; subsequent iterations use the finer mesh. Early iterations are 2--6$\times$ faster.

\paragraph{Dual Backends.}
DOLFINx (default) and scikit-fem are supported. DOLFINx is typically 2.5--3$\times$ faster per solve; scikit-fem is used when DOLFINx is unavailable or when \texttt{structural\_backend: skfem} is set. Both evaluate J$\times$B at quadrature points and use P1 tetrahedra.

\begin{table*}[t]
  \centering
  \caption{Key implementation parameters.}
  \begin{tabular}{lll}
    \toprule
    Parameter & Default & Description \\
    \midrule
    \texttt{width}, \texttt{height} & 0.05\,m & Cross-section dimensions \\
    \texttt{E} & 100\,GPa & Young's modulus \\
    \texttt{nu} & 0.3 & Poisson ratio \\
    \texttt{mesh\_resolution} & 0.16\,m & Gmsh element size \\
    \texttt{stress\_metric} & $\max_e \sigma_{\text{vm}}^{(e)}$ & Scalar metric \\
    \texttt{fd\_step} & $10^{-5}$ & Finite-difference step \\
    \texttt{eval\_interval} & 1 & Evaluate every $N$ iterations \\
    \texttt{BC\_Z\_FRACTION} & 0.15 & Fraction of $z$-range pinned \\
    \texttt{structural\_quadrature\_degree} & 2 & Quadrature degree for J$\times$B (1=centroid, 2=full) \\
    \texttt{structural\_mesh\_resolution\_coarse} & 0.16\,m & Coarse mesh for adaptive strategy \\
    \texttt{structural\_mesh\_resolution\_fine} & 0.08\,m & Fine mesh after refinement \\
    \texttt{structural\_refine\_stress\_ratio} & 0.5 & Refine when $\sigma_{\text{vm}} \geq \eta \times$ threshold \\
    \texttt{structural\_use\_cached\_K} & false & Cache stiffness $\mathbf{K}$ for $\nabla J$ speedup \\
    \texttt{structural\_backend} & dolfinx & FEM backend: dolfinx or skfem \\
    \bottomrule
  \end{tabular}
  \label{tab:fem_summary}
\end{table*}
 or \include{...}

\section{Methodology for in-the-loop FEM}\label{sec:appendix_B}

FEM calculations must be optimized for speed in order to be computed in the optimization loop. We detail below a number of choices that were made to implement, make efficient, and make robust the in-the-loop FEM calculations presented in this work.

\paragraph{Domain and Geometry.}
The coil winding pack is represented as a solid with rectangular cross-section swept along the coil centerline. The mesh is a tetrahedralization of this volume, generated either by Gmsh or a pre-built mesh from ParaStell or Gmsh.

\paragraph{Mesh Deformation Strategy.}
The mesh topology is generated once and cached. On subsequent evaluations, node positions are deformed to follow the updated coil geometry:
\begin{enumerate}
  \item For each mesh node $\mathbf{x}$, find the nearest coil centerline point $\boldsymbol{\gamma}_{i}$ and the associated coil frame $\{\mathbf{p}_i, \mathbf{q}_i, \mathbf{t}_i\}$.
  \item Compute the local offset in the coil frame:
    \begin{equation}
      \delta_p = (\mathbf{x} - \boldsymbol{\gamma}_i) \cdot \mathbf{p}_i,\quad
      \delta_q = (\mathbf{x} - \boldsymbol{\gamma}_i) \cdot \mathbf{q}_i,\quad
      \delta_t = (\mathbf{x} - \boldsymbol{\gamma}_i) \cdot \mathbf{t}_i.
    \end{equation}
  \item When the coil deforms, reconstruct the new node position:
    \begin{equation}
      \mathbf{x}_{\text{new}} = \boldsymbol{\gamma}_i + \delta_p \, \mathbf{p}_i + \delta_q \, \mathbf{q}_i + \delta_t \, \mathbf{t}_i.
    \end{equation}
\end{enumerate}

The coil frame uses the centroid frame: $\mathbf{C} = \operatorname{mean}(\boldsymbol{\gamma})$, $\mathbf{w} = \boldsymbol{\gamma} - \mathbf{C}$, $\mathbf{p} = (\mathbf{w} - (\mathbf{w} \cdot \mathbf{t})\mathbf{t}) / |\cdots|$, $\mathbf{q} = \mathbf{t} \times \mathbf{p}$.

\paragraph{Current Density Model.}
The body force density is $\bm f = \bm J \times \bm B$, which requires a model for the current in the coil and the magnetic field and must be evaluated on the appropriate quadrature points of the weak formulation described below.
The current density is assumed uniform across the winding-pack cross-section:
\begin{equation}
  \mathbf{J} = \frac{I}{A} \, \mathbf{t},
\end{equation}
where $I$ is the coil current, $A$ is the cross-section area, and $\mathbf{t}$ is the unit tangent at the nearest centerline point. Each mesh node (or quadrature point) is assigned to the nearest coil centerline point via a KD-tree lookup.

\paragraph{Magnetic Field model.}
The Biot--Savart law diverges near the coil centerline, so a filamentary model is inadequate for points inside the conductor. StellCoilBench uses the Landreman et al.\ (2025) regularized internal field model for rectangular cross-sections. The total field at a point inside a coil is:
\begin{equation}
  \mathbf{B} = \mathbf{B}_{\text{reg}} + \mathbf{B}_0 + \mathbf{B}_{\kappa} + \mathbf{B}_b.
\end{equation}
This is a fast approximation for the self-field that significantly increases the capacity to compute the Lorentz body force in-the-loop.
The mutual field is computed via the Biot-Savart law assuming filamentary models for the other coils, since this is adequate outside the coil where the calculations are taking place (assuming other coils do not come too close).

\paragraph{Linear Elasticity.}
The linear structural problem to solve is: find the displacement $\mathbf{u} : \Omega \to \mathbb{R}^3$ such that
\begin{equation}
  -\nabla \cdot \boldsymbol{\sigma} = \mathbf{f}
  \quad \text{in } \Omega,
  \qquad
  \mathbf{u} = \mathbf{0}
  \quad \text{on } \Gamma_\Omega,
\end{equation}
where $\boldsymbol{\sigma}$ is the Cauchy stress tensor for a linear elastic, isotropic material:
\begin{equation}
  \boldsymbol{\sigma} = \lambda \, (\nabla \cdot \mathbf{u}) \, \mathbf{I} + 2\mu \, \boldsymbol{\varepsilon},
  \qquad
  \boldsymbol{\varepsilon} = \frac{1}{2}\bigl(\nabla \mathbf{u} + (\nabla \mathbf{u})^T\bigr).
\end{equation}
The Lam\'e Parameters are:
\begin{equation}
  \lambda = \frac{E \nu}{(1+\nu)(1-2\nu)},
  \qquad
  \mu = \frac{E}{2(1+\nu)},
\end{equation}
where $E$ is Young's modulus [Pa] and $\nu$ is Poisson ratio. The defaults are $E = 100\,\text{GPa}$, $\nu = 0.3$, which are typical for REBCO and similar materials~\cite{srivastava2024modelling}.
Simple ``fixed support'' boundary conditions are used; nodes whose $z$-coordinate lies in the bottom fraction $\zeta = 0.15$ of the global mesh $z$-range are pinned ($\mathbf{u} = \mathbf{0}$). This models coil support stands and removes rigid-body modes:
\begin{equation}
  \Gamma_\Omega = \bigl\{ \mathbf{x} \in \partial\Omega : z(\mathbf{x}) \leq z_{\min} + 0.15(z_{\max} - z_{\min}) \bigr\}.
\end{equation}
More complex boundary conditions consistent with realistic support structures can easily be implemented in the code in future work.

\paragraph{Weak Form and Discretization.}
Using linear tetrahedral elements and a standard Galerkin discretization:
\begin{equation}
  \int_\Omega \bigl( \lambda \, (\nabla \cdot \mathbf{u}) \, (\nabla \cdot \mathbf{v}) + 2\mu \, \boldsymbol{\varepsilon}(\mathbf{u}) : \boldsymbol{\varepsilon}(\mathbf{v}) \bigr) \, dV
  = \int_\Omega \mathbf{f} \cdot \mathbf{v} \, dV
  \quad \forall \mathbf{v} \in V_h.
\end{equation}
The quadrature degree for the load integral is configurable (see Reduced Quadrature below).
The implementation uses scikit-fem or DOLFINx with bilinear and linear forms assembled and the linear system solved with Dirichlet DOFs eliminated via condense.

\paragraph{Von Mises Stress.}
For each tetrahedral element, the displacement gradient $\nabla \mathbf{u}$ is computed from nodal values and the element Jacobian. The strain tensor is
\begin{equation}
  \boldsymbol{\varepsilon} = \frac{1}{2}\bigl(\nabla \mathbf{u} + (\nabla \mathbf{u})^T\bigr).
\end{equation}
The Cauchy stress is
\begin{equation}
  \boldsymbol{\sigma} = \lambda \, \operatorname{tr}(\boldsymbol{\varepsilon}) \, \mathbf{I} + 2\mu \, \boldsymbol{\varepsilon}.
\end{equation}
The deviatoric stress is $\mathbf{s} = \boldsymbol{\sigma} - \frac{1}{3}\operatorname{tr}(\boldsymbol{\sigma}) \, \mathbf{I}$. The Von Mises stress is
\begin{equation}
  \sigma_{\text{vm}}
  = \sqrt{\frac{3}{2} \sum_{i,j} s_{ij}^2}.
\end{equation}
This is computed per element (constant within each linear tetrahedron).
The FEM yields $\sigma_{\text{vm}}^{(e)}$ per element $e$, so that the objective is chosen to be:
\begin{equation}
  \max(0, \max_e \sigma_{\text{vm}}^{(e)} - \lambda_{\sigma_{vm}})^2.
\end{equation}
The objective returns values in GPa to improve optimizer numerics.

\paragraph{Optimization tricks.}
We use a number of optimization tricks to prevent undesirable behavior.
The first is a guard against computing FEM when coil-coil distances become very small, since the mutual field body force calculations will become nearly singular. When $d_{cc} < \alpha_\text{safety}\lambda_{d_{cc}}$ it skips the FEM and returns a finite penalty (e.g., 10\,GPa) with zero gradient.
The second trick is a short-circuit; when $\max_e \sigma_{\text{vm}}^{(e)} \leq \lambda_{\sigma_{vm}}$, it skips the expensive gradient calculation and returns $\nabla J = \mathbf{0}$.
Gradients are computed by forward finite differences.
Third, the stiffness matrix $\mathbf{K}$ and fixed-DOF list can be cached at the baseline; $\mathbf{K}$ changes by $O(\varepsilon)$ per perturbation, which is negligible for FD accuracy. When \texttt{structural\_use\_cached\_K} is true, $\mathbf{K}$ is reused for all FD perturbations in $\nabla J$, yielding up to $\sim$40\,\% speedup on $\nabla J$ for adaptive meshes at the cost of gradient accuracy; the default is false.
Fourth, only coil frames and $\mathbf{B}_{\text{reg}}$ for coils whose DOFs are perturbed are recomputed, reducing cost.
Fifth, the objective can be evaluated every $N$ iterations and cached in between to reduce cost.
Both DOLFINx and scikit-fem evaluate J$\times$B at quadrature points; the backend choice affects per-solve speed (see Dual Backends below).

\paragraph{MPI Parallelization.}
When running with MPI ($P > 1$ processes), the structural gradient $\nabla J$ is computed in parallel: rank~0 broadcasts the baseline state; each rank computes a subset of the $n$ finite-difference perturbations (indices $r, r+P, r+2P, \ldots$); partial gradients are combined via \texttt{Allreduce}. Workers (ranks $1 \ldots P{-}1$) block in a worker loop until rank~0 signals a gradient computation. Expected speedup is up to $\sim P\times$ on the FD loop; overhead from broadcast and Allreduce is small for typical $n \sim 100$--200. Example: \texttt{mpirun -n 4 stellcoilbench submit-case cases/basic\_LandremanPaulQA\_structural.yaml} uses 4 processes for parallel $\nabla J$.

\paragraph{Reduced Quadrature.}
The weak-form RHS uses a quadrature element with configurable degree $q \in \{1,2\}$. For $q=1$ (centroid-only), J$\times$B evaluations are reduced, giving $\sim$30--50\,\% speedup in body-force assembly with $\sim$2\,\% relative error in $J$ for typical coils. Configurable via \texttt{structural\_quadrature\_degree}.

\paragraph{Adaptive Mesh.}
When \texttt{structural\_mesh\_resolution\_coarse} and \texttt{structural\_mesh\_resolution\_fine} are set, optimization starts with the coarse mesh and refines to the fine mesh when $\sigma_{\text{vm}} \geq \eta \times \lambda_{\sigma_{\text{vm}}}$ ($\eta = 0.5$ by default). Refinement is one-way (coarse $\to$ fine) and triggered from the optimizer callback; subsequent iterations use the finer mesh. Early iterations are 2--6$\times$ faster.

\paragraph{Dual Backends.}
DOLFINx (default) and scikit-fem are supported. DOLFINx is typically 2.5--3$\times$ faster per solve; scikit-fem is used when DOLFINx is unavailable or when \texttt{structural\_backend: skfem} is set. Both evaluate J$\times$B at quadrature points and use P1 tetrahedra.

\begin{table*}[t]
  \centering
  \caption{Key implementation parameters.}
  \begin{tabular}{lll}
    \toprule
    Parameter & Default & Description \\
    \midrule
    \texttt{width}, \texttt{height} & 0.05\,m & Cross-section dimensions \\
    \texttt{E} & 100\,GPa & Young's modulus \\
    \texttt{nu} & 0.3 & Poisson ratio \\
    \texttt{mesh\_resolution} & 0.16\,m & Gmsh element size \\
    \texttt{stress\_metric} & $\max_e \sigma_{\text{vm}}^{(e)}$ & Scalar metric \\
    \texttt{fd\_step} & $10^{-5}$ & Finite-difference step \\
    \texttt{eval\_interval} & 1 & Evaluate every $N$ iterations \\
    \texttt{BC\_Z\_FRACTION} & 0.15 & Fraction of $z$-range pinned \\
    \texttt{structural\_quadrature\_degree} & 2 & Quadrature degree for J$\times$B (1=centroid, 2=full) \\
    \texttt{structural\_mesh\_resolution\_coarse} & 0.16\,m & Coarse mesh for adaptive strategy \\
    \texttt{structural\_mesh\_resolution\_fine} & 0.08\,m & Fine mesh after refinement \\
    \texttt{structural\_refine\_stress\_ratio} & 0.5 & Refine when $\sigma_{\text{vm}} \geq \eta \times$ threshold \\
    \texttt{structural\_use\_cached\_K} & false & Cache stiffness $\mathbf{K}$ for $\nabla J$ speedup \\
    \texttt{structural\_backend} & dolfinx & FEM backend: dolfinx or skfem \\
    \bottomrule
  \end{tabular}
  \label{tab:fem_summary}
\end{table*}
